%% kapa: exam=1 round=35 year=2024 period=1 form=A subject=economics subject_ko=경제학원론 seq=3 number=43
%% source: Q-Net 감정평가사 기출문제 (cst00309) artlSeq=5213229 file=raw/1008/5213229/1교시 문제지 원본.pdf
%% conversion: kordoc 4.24.51 -> tools/gen1.py (auto); review=proofread-sonnet; flags=tikz
\begin{question}{43}
두 재화 $X$, $Y$에 대한 갑의 효용함수가 $U=X+Y+$ min $\{X,Y\}$일 때, 갑의 무차별 곡선으로 적절한 것은?
\begin{center}
\begin{tabular}{@{}l@{\hspace{1.5em}}l@{\hspace{1.5em}}l@{}}
\begin{tikzpicture}[x=0.1mm,y=-0.1mm,font=\scriptsize,line width=0.4pt]
\node[anchor=north east,font=\small] at (14,10) {①};
\draw (38,5) -- (38,308) -- (370,308);
\node at (24,6) {\textit{Y}};
\node at (380,322) {\textit{X}};
\node at (22,322) {0};
\draw[line width=1.2pt] plot[smooth] coordinates {(88,62) (97,112) (125,175) (170,225) (220,252) (260,265) (306,273)};
\end{tikzpicture} &
\begin{tikzpicture}[x=0.1mm,y=-0.1mm,font=\scriptsize,line width=0.4pt]
\node[anchor=north east,font=\small] at (14,10) {②};
\draw (38,5) -- (38,308) -- (370,308);
\node at (24,6) {\textit{Y}};
\node at (380,322) {\textit{X}};
\node at (22,322) {0};
\draw[line width=1.2pt] (88,66) -- (330,273);
\end{tikzpicture} &
\begin{tikzpicture}[x=0.1mm,y=-0.1mm,font=\scriptsize,line width=0.4pt]
\node[anchor=north east,font=\small] at (14,10) {③};
\draw (38,5) -- (38,308) -- (370,308);
\node at (24,6) {\textit{Y}};
\node at (380,322) {\textit{X}};
\node at (22,322) {0};
\draw[dotted] (38,308) -- (340,37);
\node at (240,70) {$Y=X$};
\draw[line width=1.2pt] (118,69) -- (118,239) -- (302,240);
\end{tikzpicture}\\
\begin{tikzpicture}[x=0.1mm,y=-0.1mm,font=\scriptsize,line width=0.4pt]
\node[anchor=north east,font=\small] at (14,10) {④};
\draw (38,5) -- (38,308) -- (370,308);
\node at (24,6) {\textit{Y}};
\node at (380,322) {\textit{X}};
\node at (22,322) {0};
\draw[dotted] (38,308) -- (343,53);
\node at (250,78) {$Y=X$};
\draw[line width=1.2pt] (86,43) -- (152,203) -- (315,261);
\end{tikzpicture} &
\begin{tikzpicture}[x=0.1mm,y=-0.1mm,font=\scriptsize,line width=0.4pt]
\node[anchor=north east,font=\small] at (14,10) {⑤};
\draw (38,5) -- (38,308) -- (370,308);
\node at (24,6) {\textit{Y}};
\node at (380,322) {\textit{X}};
\node at (22,322) {0};
\draw[dotted] (38,308) -- (190,1);
\draw[dotted] (38,308) -- (318,145);
\node at (225,30) {$Y=2X$};
\node at (240,130) {$Y=\frac{1}{2}X$};
\draw[line width=1.2pt] (85,41) -- (118,108) -- (207,212) -- (332,242);
\end{tikzpicture}\\
\end{tabular}
\end{center}

\end{question}
\answer{4}
